\section{Weights and Representations}

We only treated the adjoint representation of \(\mathfrak{g}\) so far, but we can also consider other representations. In this section we will see how the roots of \(\mathfrak{g}\) are related to the weights of its representations, and how the geometry of the roots can be used to understand the structure of the representations.

We will study the eigenvalues of the Cartan subalgebra \(\mathfrak{h}\) operator in other representations.

    [...]

Thus we are basically going to repeat the same analysis we did for the adjoint representation with the roots, but now for a generic representation \((d,V)\) of \(\mathfrak{g}\), getting to the weights. We will see that the structure of the weights of \((d,V)\) is closely related to the structure of the roots of \(\mathfrak{g}\), and that we can use the geometry of the roots to understand the structure of the weights.

\begin{definition}
    Given a complex representation \((d,V)\) of \(\mathfrak{g}\) with a fixed Cartan subalgebra \(\mathfrak{h}\), an element \(\omega \in \mathfrak{h}^\vee\) of the vector space dual to the Cartan subalgebra is called a weight of \((d,V)\), if
    \[
        \dots
    \]
    Conversely...
    \[
        \dots
    \]
    [...]
    \[
        \Omega\dots
    \]

\end{definition}
This is the analogous construction to the roots, but now for a generic representation instead of the adjoint representation, defining the \textbf{weight system}.

By standard diagonalization theorems, this implies a \textbf{weight decomposition} of
\[
    V = \bigoplus_{\omega \in \Omega_d} V_\omega,
\]
of the representation space, similarly to the root  ecomposition of \(\mathfrak{g}\). For a finite dimensional representation, it also means that the set \(\Omega_d\) of weights is finite.

\begin{example}[\(\mathfrak{g} = \mathfrak{sl}(2, \mathbb{C})\)]
    For \(\mathfrak{g}= \mathfrak{sl}(2, \mathbb{C})\), the irreps of spin \(j\) have the obvious decomposition
        [...]
    \[
        d(J_3) \ket{j,m} = m \ket{j,m}
    \]
    ...
\end{example}

By diagonalizing the cartan subalgebra, we immediatly know how the ... act on the weight spaces, and we can use this to understand the structure of the representation.
\[
    \mathfrak{g} = \mathfrak{h} \oplus \bigoplus \mathfrak{g}_{\alpha},
\]
where...
By taking \(E_{\alpha} \in \mathfrak{g}_{\alpha}\) and \(y \in V_\omega\) a weight vector, we have by direct computation that
\[
    \begin{aligned}
        \dots
    \end{aligned}
\]
So the Lie algebra elements in \(\mathfrak{g}_{\alpha}\) raise / lower weights by the positive / negative root \(\alpha\). Again, this is familiar from the \(\mathfrak{sl}(2, \mathbb{C})\) example. Starting from a weight \(\omega\), we can raise or decrease the cartan eigenvalue by adding or subtracting roots, and we can use this to understand the structure of the representation.

\begin{proposition}
    If \((d,V)\) is an irrep of \(\mathfrak{g}\), then the difference of any two weights is in the root lattice:
    \[
        \dots
    \]
    such that it is an integer linear combination of roots.
\end{proposition}
\begin{proof}
    [...]
    invariant since cartan dont change anything and the other operators from the weight spaces just add or subtract weights, so we can only get weights that differ by integer linear combinations of weights.
\end{proof}

we can use the \(\mathfrak{sl}(2, \mathbb{C})\) ...
\[
    \dots
\]
where now we think of \(\omega \in \mathfrak{h}_\mathbb{R}\).
...

\begin{theorem}
    The weights \(\omega\) of any representation of \(\mathfrak{g}\) satisfy
    \[
        \dots
    \]
    the integrality condition.
\end{theorem}


[Definition and Proposition 3.6.4.]


\begin{example}[\(\mathfrak{g} = \mathfrak{su}(2)_\mathbb{C} = \mathfrak{sl}(2, \mathbb{C})\)]
    ...
    The foundamentl weight is the one sending \(J_3\) to \(\tfrac{1}{2}\).
    It is the highest weight of the spin \(\tfrac{1}{2}\) representation, and it is the only weight of that representation. The spin \(j\) representation has highest weight \(j\) and weights \(j, j-1, \ldots, -j\).

    Since the weight lattice the \(\mathbb{Z}\)-span of the fundamental weights, it means that any weight \(\omega\) is an integer multiple of \(w\)

    \dots
\end{example}

\begin{example}[\(\mathfrak{g} = \mathfrak{su}(3)_\mathbb{C} = \mathfrak{sl}(3, \mathbb{C})\)]
    We have already seen an explicit parametrization of the lie algebra \(\mathfrak{sl}(3, \mathbb{C})\) in terms of the Gell-Mann matrices, and we can use that to find the roots and weights of the adjoint representation. We have already found the two roots, after computing the commutators with the Cartan subalgebra.
    ...

    \[
        \Delta = \dots
    \]

    ...

    \[
        K(H_i,H_j) = \Tr_{\mathfrak{g}}(\mathrm{ad}_{H_i} \mathrm{ad}_{H_j}) = \kappa_{ij}
    \]
    and
    \[
        \Tr_V (d(H_i) d(H_j)) = \lambda \kappa_{ij} = C_{ij}
    \]
    where \((d,V)\) is the defining representation of \(\mathfrak{sl}(3, \mathbb{C})\) on \(\mathbb{C}^3\), and \(\lambda\) is a constant that depends on the normalization of the trace. But when computing the weights, we are only interested in the ratio of these two traces, so the constant \(\lambda\) will not affect the result.
    ...
    Since we want to find the foundamental weight, we need to find the ...

    We want to find the parametrization \(\omega = a_1 \alpha^{(1)} + a_2 \alpha^{(2)}\) such that the following is a foundamental weight:
    \[
        \left\langle \alpha^{(k)},\, \omega \right\rangle = \dots
    \]
    ...
    think about raw vector space, and then we can use the inner product to find the coordinates of the fundamental weights in the root basis. Just algebra and basis change, we can find the coordinates of the fundamental weights in the root basis, and then we can use that to find the coordinates of the roots in the fundamental weight basis, which is what we want to plot them in.
    ...

        [\textit{drawing: plusses are previous hexagon roots, fundamental weights are w1 and w2 in blue and orange, then applying linear combinations of the roots to the fundamental weights we get the other weights of the defining representation, which are the vertices of the triangle (one blue and one orange)}]

    What is the relation among this geometry and the structure of the representation?

    \[
        w^{(1)} = (1,0) \implies W^{(i)} \in \mathfrak{h}^\vee \text{ such that } \dots
    \]
    def rep and action of \(H_1\) and \(H_2\) on canonical basis column vectors, looking at the eigenvalues we find the weights of the defining representation, then he rewrites as the sum of the fundamental weights and the roots.

        [...]

    The three blue dots then represented the weights of the defining representation, which are the vertices of the triangle, and the six plusses are the roots, which are the vertices of the hexagon. Then the orange ones related to the complex conjugate of the defining representation, which is the anti-fundamental representation, and its weights are the vertices of the triangle with orange vertices. They are indeed made by subtracting the roots from the fundamental weights, ...
    The complex conjugate and the foundamental representations are not the same since they have different weights.
        [...]

\end{example}

